\section{Experimental Design}
\subsection{Experimental Objectives and Conditions}
The experiment was designed to evaluate the system-level effect of object-conditioned joint parameter configuration during heterogeneous-object haptic teleoperation. Task completion time was the primary endpoint. The primary planned comparison was between the complete joint parameter-bundle mode (C) and the impedance-only reconfiguration mode (E). These modes shared the same visual-semantic display and slave-side impedance update; mode C additionally updated the haptic-interface and gripper-execution parameter groups. Thus, the C--E comparison tests the system-level contribution of the complete bundle beyond impedance-only reconfiguration, but does not identify the independent effects of the haptic and gripper groups or an interaction between them.

Comparisons of mode C with the fixed-parameter baseline (A), the manual-selection workflow baseline (B), and the visual-semantic display with fixed parameters (D) were specified as secondary planned contrasts. Success rate and Raw NASA-TLX were secondary descriptive outcomes. Master-side trajectory length and pause count were exploratory process measures. All analyses were interpreted as repeated-measures evidence within the tested three-operator sample rather than as population-level human-factors inference.

Table~\ref{tab:3} defines the five experimental modes, including the visual information available to the operator, the parameter groups that could change, and the role of each mode in the comparison. The visual-semantic display remained active throughout formal trials in modes C--E. Modes C and D held the visual-semantic display constant while differing in joint parameter configuration. Mode B required manual object identification and button-based strategy selection after timing had begun; its comparison with mode C consequently includes workflow automation and interaction differences and is not interpreted as an isolated controller effect.

\begin{table*}[!t]
\centering
\caption{Experimental modes, channel-update logic, and comparison roles.}
\label{tab:3}
\scriptsize
\setlength{\tabcolsep}{4pt}
\renewcommand{\arraystretch}{1.20}
\begin{tabularx}{\textwidth}{@{}p{0.13\textwidth}p{0.23\textwidth}>{\centering\arraybackslash}p{0.105\textwidth}>{\centering\arraybackslash}p{0.105\textwidth}>{\centering\arraybackslash}p{0.105\textwidth}Y@{}}
\toprule
Mode & Visual cue / selection & Slave impedance & Haptic interface & Gripper execution & Comparison role \\
\midrule
A: Fixed-parameter baseline & None & --- & --- & --- & Fixed-parameter baseline. \\
B: Manual-selection workflow baseline & No visual cue; operator selects a complete strategy after timing starts. & Manual & Manual & Manual & Workflow baseline; timing includes manual identification and selection. \\
C: Joint parameter-bundle configuration & Class, strategy, confidence, and colored bounding box & Auto & Auto & Auto & Full joint configuration. \\
D: Visual-semantic display with fixed parameters & Same as C & --- & --- & --- & Visual-semantic display control. \\
E: Impedance-only reconfiguration & Same as C & Auto & --- & --- & Impedance-only control. \\
\bottomrule
\end{tabularx}
\par\smallskip
\footnotesize\textit{Note:} Auto denotes a one-shot update after the first valid visual detection; Manual denotes manual strategy selection after timing starts; --- retains the initialized setting. In modes C--E, no detection or confidence below 0.25 retains the initialized state. A detected class without an explicit mapping invokes the balanced strategy; in mode D, the visual-semantic output is displayed while $\Theta_0$ is retained. Vision is not used in modes A and B. The C--E contrast changes the haptic-interface and gripper parameter groups together and therefore does not identify their individual effects or an interaction between them.
\end{table*}

\subsection{Operators, Objects, and Matched Block Design}
Three operators completed the main experiment (P01--P03; 23--24 years old; two male and one female; all right-handed). The Omega.7 was positioned on the left side of the platform and was operated with the left hand in every trial. This arrangement was held constant across operators and experimental modes to maintain a common physical interface condition. All operators had received basic teleoperation training and completed a 10--15 min familiarization session immediately before formal data collection to practice master-device motion, gripper input, and the task procedure. All operators provided written informed consent.

The six objects were selected to span different contact, grasping, and transfer demands, including deformability, smooth or irregular surfaces, and elongated geometry. They were assigned to the fragility-priority, balanced, or stability-priority strategy solely to select the parameter settings used in this platform and task. These assignments do not constitute a general classification of object material properties. Table~\ref{tab:5} lists the physical profiles, assigned strategies, and principal manipulation risks; the contrasting paper-cup and scissors cases illustrate the balance between deformation avoidance and orientation retention.

\begin{table*}[!t]
\centering
\caption{Test objects, assigned strategies, and principal task risks.}
\label{tab:5}
\small
\setlength{\tabcolsep}{6pt}
\renewcommand{\arraystretch}{1.16}
\begin{tabularx}{\textwidth}{@{}p{0.14\textwidth}p{0.18\textwidth}p{0.29\textwidth}Y@{}}
\toprule
Object & Assigned strategy & Physical profile & Principal task risk \\
\midrule
Apple & Fragility-priority & $\sim$200 g; smooth; $\varnothing$70--80 mm & Impact or slip; gentle contact required. \\
Banana & Fragility-priority & $\sim$120 g; smooth; $20\times180$ mm & Compression-induced deformation and slip. \\
Paper cup & Balanced & $\sim$5 g; paper; $\varnothing$75$\times$90 mm & Easily deformed; stable grasp required. \\
Bottle & Balanced & $\sim$30 g; smooth plastic; $\varnothing$65$\times$200 mm & Slip; balance between efficiency and stability. \\
Mouse & Stability-priority & $\sim$100 g; smooth plastic; $65\times120\times35$ mm & Rigid irregular surface; transfer slip. \\
Scissors & Stability-priority & $\sim$150 g; metal and plastic; $50\times170\times15$ mm & Rigid elongated geometry; high orientation demand. \\
\bottomrule
\end{tabularx}
\end{table*}

The experiment comprised 27 matched task blocks. A block was defined by a fixed operator, object, and repetition index and contained one trial in each of modes A--E. Accordingly, $27$ blocks $\times$ $5$ modes per block yielded $135$ trials. This structure makes the five mode observations within each block the pairing unit for the completion-time analysis.

Within each operator and strategy category, the two objects were assigned to the three repetitions in a fixed 2:1 ratio. The allocation was held constant across the five modes within a block, preserving the object-matched comparison, and the 2:1 allocation was reversed across operators. The resulting object counts gave the approximate 5:4 balance within each strategy shown in Table~\ref{tab:6}. This allocation was specified by the experimental schedule, rather than selected after observing performance.

\begin{table}[!t]
\centering
\caption{Allocation of matched task blocks and trials by strategy and object.}
\label{tab:6}
\footnotesize
\setlength{\tabcolsep}{6pt}
\renewcommand{\arraystretch}{1.16}
\begin{tabular*}{\columnwidth}{@{\extracolsep{\fill}}lcc@{}}
\toprule
Object & Blocks & Trials \\
\midrule
\multicolumn{3}{@{}l}{\textit{Fragility-priority}} \\
Apple & 4 & 20 \\
Banana & 5 & 25 \\
\addlinespace[2pt]
\multicolumn{3}{@{}l}{\textit{Balanced}} \\
Paper cup & 5 & 25 \\
Bottle & 4 & 20 \\
\addlinespace[2pt]
\multicolumn{3}{@{}l}{\textit{Stability-priority}} \\
Mouse & 5 & 25 \\
Scissors & 4 & 20 \\
\midrule
Total (six objects) & 27 & 135 \\
\bottomrule
\end{tabular*}
\par\smallskip
\footnotesize\textit{Note:} Each block contains one trial in each of the five modes.
\end{table}

Five-mode orders were preassigned at the operator-by-strategy-group level using nine different sequences, so the design did not rely on a single fixed order. The schedule distributed modes across early and late positions but was not a complete counterbalanced Latin-square design. Across the nine sequences, mode C occurred in position four in three groups, whereas mode E occurred in position five in three groups. The exact sequences are reported in Supplementary Table S2. With three operators, order, object, and operator effects could not be fully separated; C--E is therefore interpreted as a within-sample paired contrast rather than an order-free population effect.

\subsection{Task Procedure and Outcome Measures}
Before each trial, the test object was manually placed within a predefined work region and at an approximately prescribed orientation. Small reset errors in position and orientation were unavoidable, but the same placement protocol was used for all modes. Only translational master motion was mapped; the desired Panda end-effector orientation remained fixed throughout the trial.

Each trial followed reset, approach, grasp, transfer, release, and termination stages. In modes C--E, camera acquisition, visual inference, and teleoperation started concurrently at task initiation; visual output did not gate operator motion. In mode B, timing began before object identification and strategy selection, and robot motion was permitted only after the operator completed the button selection. No separate timestamp marked the end of this manual selection step.

A trial was successful when grasping, transfer, and placement were completed without a drop, observable slip, or visible damage. Complete separation of the object from the gripper was classified as a drop; visible relative motion between object and gripper during grasp or transfer as slip; and a clear indentation, crushing deformation, or other visible damage as damage. A designated researcher applied these predefined criteria in real time. Trials were not systematically video-recorded, and outcome labels were not independently or blindly re-evaluated.

Failure did not stop the timer immediately. Regrasping after a drop was not permitted; the failure was recorded, and the operator completed the prescribed return-and-termination procedure until the robot and gripper reached the predefined terminal state. The same timing endpoint was therefore applied to successful and failed trials. The system recorded master-side trajectory, gripper input, active control parameters, task duration, and outcome. It did not retain event-level visual classes, confidences, strategy-trigger timestamps, or synchronized trigger--contact timestamps; consequently, the retained record cannot retrospectively verify pre-contact strategy selection in every trial.

Task completion time was measured from task initiation to the common terminal state described above and was retained for successful and failed trials. Success rate was calculated as the number of successful trials divided by all trials in a mode and is reported descriptively because outcome labels were assigned in real time without independent review.

Subjective workload was assessed with Raw NASA-TLX \cite{ref27}. Each of the six dimensions was rated from 0 to 100, and their unweighted arithmetic mean was reported. Each operator completed one questionnaire immediately after the three trials in a given strategy $\times$ mode condition. These trials combined the two objects assigned to that strategy in the operator-specific 2:1 ratio. The workload dataset therefore contained $3$ operators $\times 3$ strategies $\times 5$ modes $=45$ questionnaires, or nine questionnaire units per mode. It supports mode- and strategy-level descriptive comparisons, but not object-specific or population-level workload inference.

Master-side trajectory length and pause count were treated as exploratory process measures. Visual performance was characterized separately in a controlled image test using class accuracy, class-to-strategy mapping accuracy, confidence, and per-image processing time. A pause was operationally defined before analysis as a continuous interval in which master-device speed was below 0.005 m/s for at least 0.30 s. Pauses were obtained by differentiating raw master-side trajectories sampled at approximately 200 Hz.

\subsection{Statistical Analysis}
Completion time across the five modes was first assessed with a Friedman test using the 27 matched task blocks as the repeated-measures units. When the omnibus test was significant, paired Wilcoxon signed-rank tests with Holm--Bonferroni correction were used for the planned contrasts. For the primary C--E contrast, the analysis reports the paired mean difference, relative change, and a 95\% bootstrap confidence interval based on 10,000 resamples of matched task blocks. Effect direction was also summarized by operator, together with a leave-one-operator-out sensitivity analysis, to assess whether a single operator dominated the observed difference. Because blocks were nested within operators, these analyses characterize conditional repeated-measures evidence in the observed sample rather than independent population-level observations.

Paired C--E success outcomes were tabulated, and an exact McNemar test was reported only as a supplementary analysis because few discordant pairs were observed. Raw NASA-TLX was compared descriptively across the nine operator $\times$ strategy questionnaire units per mode; no bootstrap interval was calculated. Trajectory length and pause count were analyzed as exploratory measures. Results are reported, as appropriate, as median [IQR] and mean $\pm$ SD. Because the 27 blocks are nested within only three operators, all inferential results and bootstrap intervals are conditional on the present repeated-measures sample and do not support population-level inference across operators.
